% Author: Simon-Pierre Boucher — contact@spboucher.ai
%
% ═══════════════════════════════════════════════════════════════════════
% 2. LITERATURE REVIEW
% ═══════════════════════════════════════════════════════════════════════
\section{Literature Review}
\label{sec:literature}

This paper sits at the intersection of four literatures: the hedonic pricing
tradition in housing economics, the econometrics of distributional heterogeneity,
the rapidly growing machine-learning strand of property valuation, and the
methodological literature on validating predictive models under spatial dependence.
We review each in turn, emphasizing in every case how it bears on the design choices
of this study and where the gap addressed by this paper lies.

\subsection{Hedonic Pricing: Origins, Theory, and the Interpretation of Implicit Prices}

The practice of regressing prices on product characteristics predates its
theoretical justification by several decades. \citet{waugh1928quality} priced
quality attributes of vegetables, \citet{court1939hedonic} constructed hedonic
price indexes for automobiles, and \citet{griliches1961hedonic} revived the method
for quality-adjusted price measurement. In housing, \citet{ridker1967determinants}
provided the first regression-based estimates of how a local disamenity---air
pollution---is capitalized into residential property values, inaugurating the
property-value approach to valuing local public goods. The theoretical foundation
arrived with \citet{lancaster1966new}, who recast goods as bundles of
characteristics, and \citet{rosen1974hedonic}, who showed that in a competitive
market the equilibrium price schedule $P(\mathbf{z})$ traces out a double envelope
of bid and offer functions, so that the gradient $\partial P/\partial z_k$ equals
the marginal implicit price of attribute $z_k$. \citet{sirmans2005composition}
synthesize 125 empirical applications, documenting robust positive gradients for
living area, bathrooms, and garages and negative gradients for property age---the
same qualitative pattern our conditional-mean estimates reproduce at national scale.

A central and often underappreciated distinction is that first-stage hedonic
gradients are equilibrium price phenomena, not preference parameters. Recovering
willingness-to-pay functions requires a second stage whose identification demands
exclusion restrictions rarely available in practice \citep{bartik1987estimation,
ekeland2004identification}. Modern syntheses formalize when hedonic estimates can
be trusted: \citet{kuminoff2013new} survey the equilibrium-sorting framework that
now underpins structural interpretation of housing-market regressions, and
\citet{bishop2020best} distill best practices for credible hedonic estimation,
emphasizing fine spatial controls, robustness to specification, and transparent
treatment of data limitations. Particularly relevant to our fixed-effects
progression, \citet{kuminoff2010which} show in large-scale simulations that
specifications with rich spatial fixed effects dramatically outperform sparse
cross-sectional specifications in recovering marginal willingness to pay when
omitted spatial variables correlate with amenities. Our region~$\rightarrow$
state~$\rightarrow$ ZIP3 exercise provides complementary, purely predictive
evidence on the same point. Throughout, we follow this literature's counsel and
interpret coefficients as conditional associations---listing-price capitalization
gradients---rather than structural demand parameters.

\subsection{Functional Form and Specification}

The semi-log specification became the workhorse of applied hedonics after
\citet{cropper1988choice} showed in Monte Carlo experiments that simple forms
outperform flexible ones when some attributes are omitted or measured with
error---precisely the situation of scraped listing data; \citet{hill2013hedonic}
surveys the subsequent evolution of hedonic specification practice in residential
applications. Earlier work had already
cautioned against atheoretic flexibility: \citet{halvorsen1981choice} demonstrated
the sensitivity of Box-Cox hedonic estimates, and \citet{anglin1996semiparametric}
documented gains from semiparametric estimation of the price function. This
literature frames one of our central questions: how much of the predictive
advantage of gradient-boosted trees reflects genuine non-linearity that the
semi-log form misses, and how much reflects something else---in our case, spatial
memorization.

\subsection{Spatial Dependence in Housing Markets}

Housing data are spatial data. \citet{can1992specification} showed that ignoring
spatial autocorrelation biases hedonic coefficient estimates;
\citet{dubin1998spatial} provides an accessible treatment of why residential
prices are spatially correlated (shared unobserved neighborhood attributes,
market-mediated spillovers); and \citet{basu1998analysis} document strong
spatial autocorrelation in Dallas transaction prices even after conditioning on
rich attribute sets. The formal apparatus of spatial econometrics---spatial lag,
spatial error, and spatial Durbin models---is developed in \citet{anselin1988spatial}
and \citet{lesage2009introduction}, with \citet{anselin2010thirty} providing a
retrospective. We deliberately stop short of estimating spatial econometric models:
our contribution is diagnostic. Using Moran's~$I$ \citep{moran1950notes} on OLS
residuals, we quantify the spatial dependence that remains after 62 attributes and
regional controls, and we show that it is a stable, structural feature of the data
rather than a sampling artifact. \citet{pace2020examining} pursue a complementary
strategy, using trees and forests to examine what information hedonic and spatial
model residuals still contain; their finding---that machine-learning methods
extract signal from residual spatial structure---anticipates our result that
tree-based models exploit location rather than only attribute non-linearity.

\subsection{Quantile Regression and Distributional Heterogeneity}

Quantile regression \citep{koenker1978regression, koenker2001quantile} replaces the
conditional mean with a family of conditional quantiles, allowing implicit prices
to vary across the price distribution. \citet{zietz2008determinants} introduced the
approach to housing, finding that square footage and other attributes are priced
differently at different quantiles; \citet{mak2010quantile} document quantile-varying
gradients in Hong Kong; and \citet{liao2012hedonic} combine quantile regression with
spatial methods. \citet{mcmillen2008changes} shows that changes in the
\emph{distribution} of house prices over time are driven largely by changes in
coefficients rather than in characteristics---evidence that quantile-specific
pricing is economically meaningful, not a statistical curiosity. More recently,
\citet{waltl2019variation} provides a comprehensive quantile analysis of the Sydney
market, documenting systematic variation across both price segments and locations.
Our contribution to this strand is scale and formality: we estimate five quantiles
on a national sample, verify coefficient stability across ten independent
subsamples, and subject the tail differences to formal inter-quantile Wald tests,
rejecting coefficient equality for 11 of 13 key attributes.

\subsection{Listing Prices, Search, and Seller Behavior}

Because our dependent variable is an asking price, the microstructure of listing
behavior matters for interpretation. Theory and evidence establish that list prices
are strategic objects: \citet{horowitz1992role} models the list price as a
commitment device in seller search; \citet{knight2002listing} shows that
overpricing lengthens time-on-market and reduces eventual sale prices;
\citet{genesove2001loss} demonstrate that loss aversion leads sellers---especially
those facing nominal losses---to set systematically higher asking prices; and
\citet{han2016role} show that the asking price plays a directing role in buyer
search, so that its information content varies across market segments. Two
implications follow for our estimates. First, listing-price gradients need not
equal transaction-price gradients, and the wedge is likely correlated with
attributes (luxury segments and distressed properties exhibit different list-to-sale
gaps). Second, this wedge affects all three modeling frameworks identically, so
\emph{comparisons across frameworks}---our primary object of interest---are
unaffected even where levels must be interpreted cautiously.

\subsection{Capitalization of Local Public Goods and Amenities}

Several of our regressors proxy local public goods, for which a mature
capitalization literature exists. \citet{oates1969effects} initiated the study of
property-tax and public-spending capitalization; \citet{black1999better} used
school-attendance boundaries to isolate the value of school quality, finding that
parents pay approximately 2\% more per 5\% increase in test scores; and
\citet{bayer2007unified} embed boundary discontinuities in an equilibrium sorting
model, showing that naive cross-sectional estimates confound school quality with
neighbor characteristics. This literature disciplines our reading of the
counterintuitive negative school-rating coefficient in the OLS results: without
boundary-style identification, school ratings in a national cross-section absorb
correlated neighborhood and fiscal variation (the property-tax rate enters
separately), and the coefficient should not be read as the value of school quality.
Similarly, \citet{pivo2011walkability} document a walkability premium using Walk
Score---the same measure we employ---while our specification, which conditions
simultaneously on bike and transit accessibility, illustrates how collinear
accessibility measures split the premium in ways that resist attribute-by-attribute
interpretation. Finally, the growing climate-capitalization literature
\citep{baldauf2020does, bernstein2019disaster, murfin2020risk} identifies a class
of price-relevant risk variables that are entirely missing from our data---a
limitation we return to in Section~\ref{sec:limitations}.

\subsection{Machine Learning in Property Valuation}

The econometrics profession has converged on a division of labor in which machine
learning excels at prediction ($\hat{y}$) problems while classical methods target
parameter ($\hat{\beta}$) problems \citep{mullainathan2017machine, varian2014big,
athey2019machine}. In real estate, this maps onto automated valuation models
(AVMs): \citet{bourassa2010predicting} compare methods for exploiting spatial
dependence in prediction; \citet{park2015using} document early machine-learning
gains in county-level housing data; \citet{kok2017big} describe the shift from
manual appraisal to big-data valuation; and \citet{steurer2021metrics} catalogue
the metrics appropriate for evaluating AVM performance, emphasizing that headline
$R^2$ figures conceal economically relevant tail behavior. The tree-ensemble
methods we deploy---random forests \citep{breiman2001random}, gradient boosting
\citep{friedman2001greedy}, and their modern implementations XGBoost
\citep{chen2016xgboost} and LightGBM \citep{ke2017lightgbm}---dominate tabular
prediction tasks of this kind. Against this backdrop, our contribution is to ask
not \emph{whether} boosting beats OLS (it does, by 20 percentage points of $R^2$
under random validation) but \emph{what that gap is made of}: the ablation design
decomposes it into feature-group contributions, and the geographic holdout reveals
that a substantial share reflects spatial memorization rather than transferable
attribute-price structure.

\subsection{Explainability: SHAP and Its Limits}

SHAP values \citep{lundberg2017unified}, rooted in the cooperative-game solution
concept of \citet{shapley1953value} and computable efficiently for trees via
TreeSHAP \citep{lundberg2020local}, have become the de facto standard for
interpreting ensemble predictions. Yet the explainability literature itself urges
caution: local surrogate explanations can be unstable \citep{ribeiro2016should},
and \citet{rudin2019stop} argues that post-hoc explanations of black-box models
should not be conflated with intrinsically interpretable modeling, particularly in
high-stakes settings. In the hedonic context the danger is specific: SHAP values
are prediction decompositions, sensitive to feature correlation, and do not satisfy
the equilibrium conditions under which Rosen's gradient equals a marginal implicit
price \citep{rosen1974hedonic, bishop2020best}. We therefore use SHAP for what it
can do---describe the predictive structure of the fitted model---and we probe the
robustness of that description by comparing importance rankings across three
different tree ensembles, finding rank correlations of 0.89--0.99.

\subsection{Validation Under Spatial Dependence}

A methodological literature largely developed in ecology and geostatistics warns
that random cross-validation overstates predictive skill whenever observations are
spatially dependent, because information leaks from training to test folds through
spatial proximity. \citet{roberts2017cross} systematize blocking strategies for
structured data; \citet{valavi2019blockcv} provide the standard software
implementation of spatial blocking; \citet{ploton2020spatial} demonstrate,
strikingly, that large-scale ecological mapping models with excellent random-CV
scores lose most of their skill under spatial validation; and
\citet{meyer2021predicting} formalize the ``area of applicability'' of spatial
prediction models. The lesson is not uncontested: \citet{wadoux2021spatial} show
that when the goal is map accuracy over a sampled region---an interpolation
problem---spatial cross-validation can be \emph{pessimistically} biased, and
design-based random validation is appropriate. This debate sharpens rather than
undermines our design: predicting prices in entirely unobserved states is an
\emph{extrapolation} task, for which held-out-region validation is the relevant
benchmark, while our random split answers the interpolation question. Reporting
both, and showing that feature sets rank differently under each, is precisely what
the debate prescribes. To our knowledge, this framing has not previously been
brought to bear on national-scale hedonic housing models.

\subsection{Research Gap}

Three gaps emerge from this review. First, the hedonic and AVM literatures rarely
meet: machine-learning papers report accuracy without engaging identification and
interpretation, while econometric papers report coefficients without assessing
predictive generalization. Comparative studies on a single large dataset that treat
OLS, quantile regression, and boosting as complementary lenses---each answering a
different question---remain scarce. Second, the spatial-validation insights of
ecology have barely penetrated housing economics, despite housing being a
canonically spatial asset; the interaction between geographic \emph{features} and
validation \emph{design} (our ablation and coordinate experiments) is essentially
unexplored at national scale. Third, the literature offers little guidance on how
explanation tools like SHAP behave across model families in housing applications.
This paper addresses all three, at the scale of 788,842 listings spanning every
U.S. state, while maintaining the interpretive discipline the identification
literature demands.
